% ID: SEL-D0083-OB-E02
% ESTADO: cotejada. Segunda pasada visual y borrador compilado revisados.
% Fuente DOC: Selectividad/sel_2012_matematicas_aplicadas/5 MATEMÁT. APLIC. CC. SS/EXAMEN 3-ANDALUCÍA.doc
% SHA-256 original: c8688a19515dfeb348e8f8817d3f4a6c0a967723c270f2ec42e7f365d001e555
% Página convertida: 2; opción B; ejercicio 2.

En el mar hay una mancha producida por una erupción submarina. La superficie afectada, en $km^2$, viene dada por la función $f(t)=\dfrac{11t+20}{t+2}$, siendo $t$ el tiempo transcurrido desde que empezamos a observarla.
\begin{enumerate}[label=\alph*)]
\item \textbf{(0.5 puntos)} \textquestiondown{}Cuál es la superficie afectada inicialmente, cuando empezamos a medirla?
\item \textbf{(1.25 puntos)} Estudie si la mancha crece o decrece con el tiempo.
\item \textbf{(0.75 puntos)} \textquestiondown{}Tiene algún límite la extensión de la superficie de la mancha?
\end{enumerate}
